% 論文節の日本語訳（読み物 PDF）用 preamble
% lualatex + jlreq
\documentclass[12pt,paper=a4]{jlreq}
\usepackage{luatexja-fontspec}
\usepackage{amsmath,amssymb,bm}
\usepackage{graphicx}
\usepackage{hyperref}
\usepackage{enumitem}
\usepackage{booktabs}
\usepackage{microtype}
\usepackage{xcolor}
\usepackage{caption}
\usepackage{tcolorbox}
\tcbuselibrary{breakable,skins}

\hypersetup{
  colorlinks=true,
  linkcolor=blue!60!black,
  urlcolor=blue!60!black,
  citecolor=blue!60!black,
}

\graphicspath{{images/}{./}}

\setlength{\parindent}{1em}
\setlength{\parskip}{0.55em}
\setlist[itemize]{leftmargin=1.6em,itemsep=0.25em,topsep=0.35em}
\setlist[enumerate]{leftmargin=1.8em,itemsep=0.25em,topsep=0.35em}

\captionsetup{
  font=small,
  labelfont=bf,
  skip=0.4em,
  justification=raggedright,
  singlelinecheck=false,
}

\newcommand{\vx}{\bm{x}}
\newcommand{\vu}{\bm{u}}
\newcommand{\veps}{\bm{\epsilon}}
\newcommand{\vw}{\bm{w}}

\DeclareMathOperator*{\argmax}{arg\,max}
\DeclareMathOperator*{\argmin}{arg\,min}
\DeclareMathOperator{\Tr}{Tr}
\DeclareMathOperator{\diag}{diag}

% 命題・定理ボックス（ページ跨ぎ時も切れ目を抑える）
\newtcolorbox{propositionbox}[1][]{%
  enhanced,
  breakable,
  colback=blue!4!white,
  colframe=blue!45!black,
  coltitle=white,
  fonttitle=\bfseries\small,
  title={#1},
  attach boxed title to top left={yshift=-2mm, xshift=3mm},
  boxed title style={
    colback=blue!45!black,
    colframe=blue!45!black,
    boxrule=0pt,
    arc=1pt,
  },
  boxrule=0.5pt,
  arc=1.5pt,
  left=5pt, right=5pt, top=8pt, bottom=5pt,
  before skip=0.9em,
  after skip=0.9em,
  pad at break=3mm,
}

% 図はフロートせず本文位置に置く（分断を防ぐ）
\newcommand{\translationfigure}[2]{%
  \begin{center}
    \includegraphics[width=0.88\linewidth]{#1}\\[0.35em]
    {\small #2\par}
  \end{center}
  \vspace{0.4em}%
}

\begin{document}
{\color{gray}\footnotesize 出典: §3.1\par}
\vspace{0.2em}
\section*{3.1 エネルギーベースモデル}



EBM に既に馴染みのある読者にとって、本節は簡潔な復習であると同時に、拡散モデルのスコアベースの見方への橋渡しでもある。

\subsection*{3.1.1 エネルギー関数による確率分布のモデル化}


データ点を $x \in \mathbb{R}^D$ とする。EBM（エネルギーベースモデル, Energy-Based Model）は、より尤もらしい配置に低いエネルギーを割り当てる、パラメータ $\phi$ でパラメータ化されたエネルギー関数 $E_\phi(x)$ を通じて確率密度を定義する。結果として得られる分布は次で与えられる：

\[
p_\phi(x) := \frac{\exp(-E_\phi(x))}{Z_\phi}, \qquad Z_\phi := \int_{\mathbb{R}^D} \exp(-E_\phi(x))\, dx,
\]


ここで $Z_\phi$ は正規化を保証する分配関数（partition function）と呼ばれる：

\[
\int_{\mathbb{R}^D} p_\phi(x)\, dx = 1.
\]


\translationfigure{fig-3-1.jpg}{図3.1: EBM の学習の模式図。「悪い」データ点では密度を下げ（エネルギーを上げ；赤矢印）、「良い」データ点では密度を上げる（エネルギーを下げ；緑矢印）。出典: 著者作成。}


この見方では、エネルギーが低い点ほど高い確率に対応し、谷へと転がり落ちるボールのイメージに近い。分配関数 $Z_\phi$ はすべての確率の和が 1 になることを保証し、結果としてエネルギーは相対的な値だけが意味を持つ。例えば、すべてのエネルギーに定数を加えても、分子と分母が同じ因子で掛けられるため、分布は変わらない。

さらに、分配関数 $Z_\phi$ が確率の和を 1 に制約するため、ある領域のエネルギーを下げればその確率は増え、補集合の確率はそれに応じて減る、ということが数学的に従う。したがって EBM は厳密な大域的トレードオフに従う：ある谷を深くすれば、他の谷は必然的に浅くなり、確率質量は各領域に独立に割り当てられるのではなく、空間全体にわたって再分配される。

\textbf{EBM における最尤推定の課題。} 原理的には、EBM は最尤推定で学習でき、これはデータへの適合と大域的正則化とを自然にバランスさせる（式 (1.1.2) を参照）：

\begin{align*}
\begin{aligned}
\mathcal{L}_{\mathrm{MLE}}(\phi)
&= \mathbb{E}_{p_{\mathrm{data}}(x)}\!\left[\log \frac{\exp(-E_\phi(x))}{Z_\phi}\right]
\quad (3.1.1)\\
&= -\underbrace{\mathbb{E}_{p_{\mathrm{data}}}[E_\phi(x)]}_{\text{データのエネルギーを下げる}}
- \log \underbrace{\int \exp(-E_\phi(x))\, dx}_{\text{大域的正則化}},
\end{aligned}
\end{align*}


ただし $Z_\phi = \int \exp(-E_\phi(x))\, dx$。第 1 項は実データのエネルギーを下げ、第 2 項は分配関数を通じて正規化を強制する。

しかし高次元では $\log Z_\phi$ とその勾配の計算は、モデル分布下の期待値を必要とするため扱いにくい。これが、対照的ダイバージェンス（contrastive divergence）（Hinton, 2002）のようにその項を近似する、あるいはスコアマッチング（score matching）によって完全に避ける、といった代替目的を動機づける。

以下では、まず §3.1.2 でスコア関数の概念を導入し、§3.1.3 で分配関数を迂回する扱いやすい学習目的としてのスコアマッチングを示し、§3.1.4 でスコア関数を用いた実用的なサンプリング手法としてのランジュバン動力学（Langevin dynamics）を論じる。

\subsection*{3.1.2 動機：スコアとは何か？}


$\mathbb{R}^D$ 上の密度 $p(x)$ に対し、スコア関数（score function）は対数密度の勾配である：

\[
s(x) := \nabla_x \log p(x), \qquad s : \mathbb{R}^D \to \mathbb{R}^D.
\]


直感的には、スコアはより高い確率の領域を指すベクトル場をなし、データが最も現れやすい場所への局所的な案内を与える（図 3.2 を参照）。

\translationfigure{fig-3-2.jpg}{図3.2: スコアベクトル場の模式図。スコアベクトル場 $\nabla_x \log p(x)$ は密度が増大する方向を示す。出典: 著者作成。}


\textbf{なぜ密度ではなくスコアをモデル化するのか？} スコアのモデル化には理論的・実用的な利点がある。

\begin{enumerate}
\item \textbf{正規化定数からの自由。} 多くの分布は、正規化されていない密度 $\tilde{p}(x)$（例: EBM における $\exp(-E_\phi(x))$）までしか定義されない：
\end{enumerate}


\[
p(x) = \frac{\tilde{p}(x)}{Z}, \qquad Z = \int \tilde{p}(x)\, dx.
\]


$Z$ の計算は扱いにくいが、スコアは $\tilde{p}$ のみに依存する：

\[
\nabla_x \log p(x) = \nabla_x \log \tilde{p}(x) - \nabla_x \log Z = \nabla_x \log \tilde{p}(x),\quad (3.1.2)
\]


（$\nabla_x \log Z = 0$、なぜなら $Z$ は $x$ について定数だからである。）これにより分配関数を完全に迂回できる。

\begin{enumerate}
\item \textbf{完全な表現。} スコア関数は背後の分布を完全に特徴づける。対数密度の勾配であるため、密度は（定数を除き）次で回復できる：
\end{enumerate}


\[
\log p(x) = \log p(x_0) + \int_0^1 s\!\bigl(x_0 + t(x - x_0)\bigr)^\top (x - x_0)\, dt,
\]


ここで $x_0$ は参照点であり、$\log p(x_0)$ は正規化によって固定される。したがって、スコアのモデル化は $p(x)$ 自体のモデル化と同じ表現力を持ちつつ、生成モデル化ではしばしば扱いやすい。

\subsection*{3.1.3 スコアマッチングによる EBM の学習}


EBM では密度は $p_\phi(x) = \frac{\exp(-E_\phi(x))}{Z_\phi}$ と定義される。最尤学習は一般に扱いにくい $Z_\phi$ の計算を要する。重要な観察は、モデル $p_\phi$ のスコアが $-\nabla_x E_\phi(x)$ に単純化し、$Z_\phi$ に依存しないことである（式 (3.1.2) を参照）。

スコアマッチング（Hyvärinen and Dayan, 2005）は、スコアがエネルギー関数のみに依存するという事実を利用する。正規化された確率を当てはめる代わりに、モデルスコアを（未知の）データスコアに揃えることで EBM を学習する：

\[
\mathcal{L}_{\mathrm{SM}}(\phi) = \tfrac{1}{2}\,\mathbb{E}_{p_{\mathrm{data}}(x)}\bigl\|\nabla_x \log p_\phi(x) - \nabla_x \log p_{\mathrm{data}}(x)\bigr\|_2^2.\quad (3.1.3)
\]


\translationfigure{fig-3-3.jpg}{図3.3: ランジュバンサンプリングの模式図。スコア関数 $\nabla_x \log p_\phi(x)$ を用い、式 (3.1.5) の更新（矢印）により軌跡を高密度領域へ導く。出典: 著者作成。}


データスコアにはアクセスできないが、部分積分により、エネルギーとその微分のみを含む等価な表現が得られる（詳細は命題 3.2.1 を参照）：

\[
\mathcal{L}_{\mathrm{SM}}(\phi)
= \mathbb{E}_{p_{\mathrm{data}}(x)}\Bigl[-\operatorname{Tr}\bigl(\nabla_x^2 E_\phi(x)\bigr) + \tfrac{1}{2}\|\nabla_x E_\phi(x)\|_2^2\Bigr] + C,
\]


ここで $\nabla_x^2 E_\phi(x)$ は $E_\phi$ のヘッセ行列であり、$C$ は $\phi$ に依存しない定数である。

この定式化は、分配関数を除去し、学習中のモデルからのサンプリングも避けるため魅力的である。主な欠点は 2 階微分の必要性である：必要となるのはヘッセ全体ではなくそのトレース $\operatorname{Tr}(\nabla_x^2 E_\phi)$ だけだが、この項の厳密な計算は一般に入力座標全体にわたる 2 階微分情報の集約を要し、コストは典型的に入力次元 $D$ とともに増える。この制約への対処は本章の後半で再訪する。

\subsection*{3.1.4 スコア関数を用いたランジュバンサンプリング}


エネルギー関数 $E_\phi(x)$ で定義される EBM からのサンプリングは、ランジュバン動力学で行える。まず離散時間のランジュバン更新を示し、次にその連続時間極限を確率微分方程式（SDE, Stochastic Differential Equation）として述べる。最後に、ランジュバン動力学が複雑なエネルギー地形の効率的な探索を可能にする物理的直感を論じる。

\textbf{離散時間ランジュバン動力学。} 離散時間のランジュバン更新は

\[
x_{n+1} = x_n - \eta\nabla_x E_\phi(x_n) + \sqrt{2\eta}\,\epsilon_n, \qquad n = 0,1,2,\ldots,\quad (3.1.4)
\]


である。ここで $x_0$ はある分布（しばしばガウス）から初期化され、$\eta > 0$ はステップサイズ、$\epsilon_n \sim \mathcal{N}(0, I)$ はガウスノイズである。ノイズは確率性を加えることで局所最小を超えた探索を可能にする。

スコア関数は

\[
\nabla_x \log p_\phi(x) = -\nabla_x E_\phi(x)
\]


と計算できるので、更新は等価に

\[
x_{n+1} = x_n + \eta\nabla_x \log p_\phi(x_n) + \sqrt{2\eta}\,\epsilon_n,\quad (3.1.5)
\]


と書ける。スコア関数がサンプルを高密度領域へ導く。この定式化は拡散モデルの中心であり、後で詳述する。

\textbf{連続時間ランジュバン動力学。} ステップサイズ $\eta$ が 0 に近づくにつれ、離散ランジュバン更新は自然に、ランジュバン確率微分方程式（Langevin SDE）で記述される連続時間過程に収束する\footnote{因子 $\sqrt{2}$ により、ランジュバン動力学は $p_\phi$ を時間について不変に保つ。すなわち $p_\phi$ は定常である：$x(0) \sim p_\phi$ ならばすべての $t \ge 0$ で $x(t) \sim p_\phi$。等価に、$p_\phi$ はフォッカー–プランク方程式（Fokker–Planck equation）（付録 B を参照）の定常解である：$\partial_t \rho = -\nabla\cdot(\rho\nabla\log p_\phi) + \frac{\sigma^2}{2}\Delta\rho$。$\rho = p_\phi$ とおくと $(\frac{\sigma^2}{2}-1)\Delta p_\phi = 0$ となり、これは $\sigma = \sqrt{2}$ のときのみ成り立つ。}：

\[
\mathrm{d}x(t) = \nabla_x \log p_\phi\bigl(x(t)\bigr)\,\mathrm{d}t + \sqrt{2}\,\mathrm{d}w(t),\quad (3.1.6)
\]


ここで $w(t)$ は標準ブラウン運動（Brownian motion）（ウィーナー過程（Wiener process）とも呼ばれる\footnote{ブラウン増分は $w(t+\eta)-w(t)\sim\mathcal{N}(0,\eta I)$ を満たす。したがってオイラー–丸山はステップノイズ $\sqrt{2}[w(t+\eta)-w(t)]=\sqrt{2\eta}\,\epsilon_n$（$\epsilon_n\sim\mathcal{N}(0,I)$）を用い、これが $\sqrt{\eta}$ 因子の由来である。ブラウン運動と SDE の詳細な導入は付録 A を参照。}）を表す。式 (3.1.4) の離散更新則は、この連続 SDE のオイラー–丸山（Euler–Maruyama）離散化として理解することが重要である。

標準的な正則性仮定の下（例: $p_\phi \propto e^{-E_\phi}$ で、$E_\phi$ が閉じ込め型で十分滑らか）、$x(t)$ の分布は $t \to \infty$ で $p_\phi$ に（指数的に速く）収束する。したがって式 (3.1.6) の SDE をシミュレート（解く）ことでサンプリングできる。

\textbf{なぜランジュバンサンプリングか？} ランジュバンサンプリングを理解する自然な方法は物理学のレンズを通すことである。エネルギー関数 $E_\phi(x)$ は粒子の振る舞いを形づくるポテンシャル地形を定義する。ニュートン動力学によれば、このエネルギーから導かれる力場の下での粒子の運動は常微分方程式（ODE）

\[
\mathrm{d}x(t) = -\nabla_x E_\phi\bigl(x(t)\bigr)\,\mathrm{d}t
\]


で記述され、これはエネルギー関数の局所最小へ粒子を決定論的に導く。しかしこのような決定論的動力学は局所最小に捕捉され、データ分布全体の探索を妨げうる。

この制約を克服するため、ランジュバン動力学は確率的摂動を導入し、次の SDE を得る：

\[
\mathrm{d}x(t) = -\nabla_x E_\phi\bigl(x(t)\bigr)\,\mathrm{d}t + \underbrace{\sqrt{2}\,\mathrm{d}w(t)}_{\text{注入ノイズ}},
\]


ここで $w(t)$ は標準ブラウン運動である。ノイズ項はエネルギー障壁を越えて局所最小から脱出することを可能にし、軌跡を確率過程にし、その定常分布はボルツマン分布

\[
p_\phi(x) \propto e^{-E_\phi(x)}
\]


に収束する。

この観点では、EBM はサンプルを高確率領域へ押しやる力場を学習するものと見なせる。ランジュバンサンプリングは、分配関数を明示的に計算せずにモデル分布 $p_\phi(x)$ からサンプルを生成する実用的な方法を与えるため、EBM に特に有用である。ランジュバン更新を反復適用することで、目標分布を近似するサンプルが得られる。

\textbf{ランジュバンサンプリングに内在する課題。} MCMC ベースのサンプラーとして広く用いられるランジュバン動力学は、高次元空間では深刻な制約に直面する。その効率は、ステップサイズ $\eta$、ノイズスケール、目標分布を正確に近似するのに必要な反復数の選択に強く依存する。

この非効率の核心は、悪い「混合時間」（mixing time）の問題にある：多くの孤立したモードを持つ複雑なデータ分布では、ランジュバンサンプリングは高確率領域間の遷移に極めて長い時間を要することが多い。この問題は次元が増えるにつれ著しく悪化し、収束が許容できないほど遅くなる。

サンプリングは、それぞれ異なるデータモードに対応する遠く離れた多くの谷を持つ広大で起伏の激しい地形を探索することだと想像できる。局所的な確率的更新に依るランジュバン動力学は、これらの谷の間を効率的に横断するのが難しい。結果として、分布の完全な多様性を捉え損ねることが多い。

この非効率さは、純粋なランダム探索よりも複雑なデータ多様体をより効果的にナビゲートできる、より構造化され誘導されたサンプリング手法の必要性を示唆する。


\end{document}
